\chapter{What we're building (and why)}

\section{The elevator pitch}
Relax.js is a small frontend library built around a Virtual DOM (VDOM) and a component model.
What makes this repository interesting is that it contains \emph{two} rendering strategies:

\begin{enumerate}
  \item A traditional VDOM runtime under \texttt{src/}.
  \item An experimental compiler-assisted runtime (HRBR) designed to update the DOM through precomputed \textbf{slots}.
\end{enumerate}

The goal of this book is not only to show \emph{how to use} Relax.js, but to explain \emph{how to build it}.
We'll read the real source code, map it to a mental model, and connect every important concept to tests.

This isn't a conceptual ``frameworks 101'' book.
It's a guided reverse-engineering and rebuild:

\begin{itemize}
  \item You will learn the smallest set of abstractions needed to describe UI.
  \item You will see how those abstractions are represented in memory.
  \item You will trace how a rendering pass turns them into real DOM.
  \item You will follow how updates are scheduled and applied.
  \item You will use the test suite as the specification that locks behavior in place.
\end{itemize}

By the end, you should be able to (1) understand every important subsystem in this repository, and (2) rebuild something similar with confidence.

\section{A tiny contract: what does Relax.js promise?}
At a high level, Relax.js promises:

\begin{itemize}
  \item You describe UI as a tree (VNode tree or compiled block).
  \item The runtime mounts it into a DOM host element.
  \item When state changes, only the necessary DOM updates happen.
  \item Components have lifecycle hooks and can emit/handle events.
\end{itemize}

Those promises are verified by tests under \texttt{src/__tests__/} and (for the HRBR runtime) under \texttt{runtime/__tests__/}.

\subsection{A quick terminology primer}
You'll see these words throughout the code and tests:

\begin{itemize}
  \item \textbf{VNode}: a plain object that describes what DOM you want.
  \item \textbf{Mount}: create DOM from a VNode for the first time.
  \item \textbf{Patch}: update existing DOM by comparing old/new VNodes.
  \item \textbf{Destroy}: remove DOM nodes and detach listeners.
  \item \textbf{Host element}: the DOM element you mount the app into (often \texttt{document.body} in tests).
  \item \textbf{HRBR block}: a compiled representation: template HTML + dynamic \emph{slots}.
  \item \textbf{Slot}: a precomputed ``write location'' (text/attr/prop/class/style/event) in a template.
\end{itemize}

\section{Public API entry points (code references)}
The \emph{public} exports live in \texttt{src/index.ts}:

\begin{lstlisting}[style=relax,language=JavaScript]
export { createApp } from './app'
export { defineComponent } from './component'
export { DOM_TYPES, h, hFragment, hSlot, hString } from './h'
export { nextTick } from './scheduler'
\end{lstlisting}

This tells you the learning order:

\begin{enumerate}
  \item \texttt{createApp()} (app lifecycle)
  \item \texttt{defineComponent()} (component model)
  \item \texttt{h()/hFragment()} (VNode model)
  \item \texttt{mountDOM()/patchDOM()/destroyDOM()} (DOM mechanics)
  \item \texttt{nextTick()} (async scheduling in the VDOM runtime)
\end{enumerate}

\section{Why a second runtime? (HRBR in one paragraph)}
The VDOM approach is flexible, but it has overhead: you allocate VNodes, diff lists, and interpret the tree every update.
HRBR is the opposite trade-off:

\begin{itemize}
  \item The compiler extracts a stable template (static HTML) once.
  \item Dynamic values become \textbf{slots} with stored DOM paths.
  \item Updates become \emph{direct writes} to known nodes.
\end{itemize}

You'll see HRBR show up in VDOM code via \texttt{DOM_TYPES.HRBR} and \texttt{hBlock()} in \texttt{src/h.ts}.

\subsection{What problem does HRBR try to solve?}
If a UI subtree is structurally stable (same elements in the same shape, only values change), a classic VDOM update does extra work:

\begin{itemize}
  \item allocate and traverse VNodes on every update
  \item interpret the tree to find where to write
  \item diff children arrays (sometimes expensive)
\end{itemize}

HRBR moves some of that cost earlier (``compile time''):

\begin{itemize}
  \item build a static DOM template once
  \item record paths to dynamic locations (slots)
  \item update by writing directly to those DOM nodes
\end{itemize}

That doesn't make VDOM obsolete.
It creates a spectrum:
VDOM is the flexible general case; HRBR is an optimized fast path when the structure is known.

\section{How we'll prove correctness throughout this book}
This repository uses \textbf{Vitest} and \textbf{jsdom}.
Our rule is simple: every time we explain a behavior, we should be able to point to a unit test.

Examples:

\begin{itemize}
  \item Mounting basics: \texttt{src/__tests__/mount-dom.test.ts}
  \item Patching/diffing: \texttt{src/__tests__/patch-dom.test.ts}
  \item Components + lifecycle: \texttt{src/__tests__/component*.test.ts}
  \item HRBR block runtime: \texttt{runtime/__tests__/block.test.ts}
\end{itemize}

\subsection{What ``tests as spec'' means in this repository}
You should treat tests here as \emph{contracts}.
When you change runtime behavior, you don't just ``fix failing tests''---you decide whether the contract changed.

Most tests follow the same simple pattern:

\begin{enumerate}
  \item Set up a small DOM (usually \texttt{document.body}).
  \item Build a VNode tree (or HRBR block).
  \item Mount, patch, or destroy.
  \item Assert on resulting \texttt{innerHTML} or element identity.
\end{enumerate}

That gives you two superpowers as a reader:

\begin{itemize}
  \item You can infer semantics without reading every implementation detail.
  \item You can reproduce bugs by writing a tiny test first.
\end{itemize}

\section{Walkthrough: the smallest end-to-end loop (with real code)}
Let's walk through the smallest app lifecycle you can observe in one sitting.

\subsection{createApp(): mount and unmount}
The application API is defined in \texttt{src/app.ts}.
The critical part is that \texttt{createApp()} remembers three pieces of state:

\begin{itemize}
  \item \texttt{parentEl}: where we mounted
  \item \texttt{isMounted}: guard against double mount/unmount
  \item \texttt{vdom}: the root VNode tree
\end{itemize}

When you mount:
\begin{enumerate}
  \item it throws if already mounted
  \item it creates the root VNode by calling \texttt{h(RootComponent, props)}
  \item it calls \texttt{mountDOM(vdom, parentEl)}
\end{enumerate}

When you unmount:
\begin{enumerate}
  \item it throws if not mounted
  \item it calls \texttt{destroyDOM(vdom)}
  \item it resets internal references
\end{enumerate}

This is a small but important architectural choice: \emph{the App owns the root VNode}.
That makes ``unmount'' a reliable cleanup boundary.

\subsection{The test that proves it works}
Open \texttt{src/__tests__/app.test.ts}.
It reads like a story:

\begin{itemize}
  \item It creates an app via \texttt{createApp(App, \{ todos: [...] \})}.
  \item It mounts to \texttt{document.body}.
  \item It waits for \texttt{nextTick()} to flush async work.
  \item It asserts a stable HTML snapshot via \texttt{singleHtmlLine\`...\`}.
\end{itemize}

Notice the discipline here: the tests make DOM updates predictable by explicitly awaiting \texttt{nextTick()}.
That foreshadows Chapter~10, where we explain why scheduling exists and how it interacts with DOM patching.

\section{Walkthrough: patching is ``same element, new values''}
Relax's VDOM runtime is built around a simple principle:

\begin{quote}
If the old and new nodes are the ``same kind'' of node, we reuse the existing DOM element and update attributes/children.
If they are different, we replace.
\end{quote}

The test suite encodes this principle over and over.
For example, \texttt{src/__tests__/patch-dom.test.ts} starts with extremely small cases:

\begin{itemize}
  \item \textbf{no change}: patching identical trees should reuse \texttt{el}.
  \item \textbf{change the root node}: patching \texttt{div -> span} replaces the root element.
  \item \textbf{patch text}: \texttt{"foo" -> "bar"} updates only the Text node.
\end{itemize}

Those tests tell you what to look for when reading \texttt{src/patch-dom.ts}: element identity reuse is a correctness requirement, not an optimization.

\section{Mini-exercise (with intent)}
Before moving on, do this intentionally (don't just skim):

\begin{enumerate}
  \item Read \texttt{src/app.ts} and answer: what runtime calls does \texttt{createApp()} delegate to?
  \item Read \texttt{src/__tests__/app.test.ts} and find every \texttt{await nextTick()}.
    For each one, explain what DOM change it is waiting for.
  \item Read the first ~5 tests in \texttt{src/__tests__/patch-dom.test.ts}.
    For each, write a one-sentence rule (e.g., ``patching text updates textContent without replacing parent elements'').
\end{enumerate}

\section{Online resources (optional, but good context)}
This book is grounded in the Relax.js repository, but it helps to know the broader ecosystem it draws from.
These are references you can read alongside chapters; they are not prerequisites.

\begin{itemize}
  \item Virtual DOM background (high-level):
  \href{https://react.dev/learn/render-and-commit}{React --- Render and Commit}.
  \item The diffing strategy commonly used in VDOM children reconciliation:
    \href{https://neil.fraser.name/writing/diff/}{Neil Fraser, ``An O(ND) Difference Algorithm and Its Variations''}.
  \item Fine-grained reactivity conceptual background (signals):
    \href{https://www.solidjs.com/docs/latest#signals}{SolidJS docs: Signals}.
  \item Scheduling and keeping UIs responsive:
    \href{https://developer.mozilla.org/en-US/docs/Web/API/window/requestAnimationFrame}{MDN: requestAnimationFrame},
    \href{https://developer.mozilla.org/en-US/docs/Web/API/MessageChannel}{MDN: MessageChannel}.
  \item Testing DOM code in Node:
    \href{https://vitest.dev/guide/}{Vitest Guide},
    \href{https://github.com/jsdom/jsdom}{jsdom}. 
\end{itemize}

If you'd like, we can make these ``real'' citations later by adding a bibliography system (BibTeX or biblatex) to \texttt{book/main.tex}.

% (Mini-exercise moved earlier and expanded.)
